\section{Future Work and Limitations}
\label{sec:future_work}

\paragraph{\cam{Scope and real-world applicability.}}
\cam{Our experiments establish the compiler's behaviour where answers can be checked: synthetic graphs with exact or MCMC ground truth (up to $n{=}8$ nodes, $64$ dimensions), plus MNIST and Colored-MNIST, where the image block is high-dimensional but the graph has two or three nodes. None of these is a real-world system, and we make no claim about performance on one. Real applications add three difficulties that our benchmarks deliberately exclude: the graph must be estimated rather than given (next paragraph), causal sufficiency may fail through hidden confounding, and graphs with many high-dimensional nodes stress both training and, more severely, evaluation (``Scaling to Large PGMs'' below). Applying StructFlow to, e.g., simulators in the sciences with known mechanistic structure, where ground truth is partially available, is the natural next test.}

\paragraph{Known structure.}
Structure discovery is out of scope: the compiler consumes a given DAG, and Appendix~\ref{app:hierarchy} delimits what is provably recovered under it---every interventional distribution for arbitrary mechanisms (Theorem~\ref{thm:level2}), exact counterfactuals for invertible-in-noise or additive mechanisms (Theorem~\ref{thm:level3})---and what is not: the graph itself is unrecoverable from observational data in the linear--Gaussian case (Proposition~\ref{prop:mec}), while for linear non-Gaussian (LiNGAM) models all three levels of Pearl's hierarchy follow end-to-end (Corollary~\ref{cor:lingam}). Coupling StructFlow to a structure learner---and studying degradation under a mis-specified graph---is the highest-value extension.

\paragraph{Counterfactual densities under partial observation.}
Corollary~\ref{cor:interventional-density} gives the interventional \emph{density} in a single pass. The counterfactual analogue only exists under \emph{partial} observation (fully-observed abduction leaves a point mass), and evaluating it requires the inference-aligned artifact alongside the causal one (Remark~\ref{rem:density-scope}); making the two artifacts cooperate on this query is a natural next step.

\paragraph{Scaling to Large PGMs.}
We evaluate exhaustively on small PGMs; the binding constraint is the cost of \emph{ground truth}, not training. The two compilations of a dense $n$-node graph jointly support ${\sim}n^2$ tractable queries (already $72$ at $n{=}8$), and every inference-direction query needs its own long NUTS reference chain at every evaluated conditioning value---this evaluation cost, not model capacity, stops our benchmarks at $n = 8$; how performance scales with graph size remains open.

\paragraph{Backbone Comparison Scope.}
The backbone ablation (Appendix~\ref{app:additional-experiments}) swaps the per-block primitive under matched budgets, isolating the compiled structure's contribution; it is \emph{not} a benchmark against modern diffusion practice---samplers, schedules, distillation, and guidance would all narrow the low-NFE gap we measure, and CNFs have accelerations we equally omit; conclusions about backbones concern the un-augmented primitives only.

\paragraph{Repairing Dynamic Conditioning at Depth.}
\cam{Section~\ref{sec:chain-stiffness} separates two errors of dynamic conditioning. The first, the bias of the block-triangular optimum, is removed by the per-block coupling. The second, a learning error that grows with depth while the coupling alone would allow far less, remains open. Natural next steps are to determine how it scales with model capacity and training length, to use couplings closer to the exact Knothe--Rosenblatt map than a batch-$128$ matching, and to use hybrid schedules that switch from dynamic to static conditioning as $t \to 1$.}
